\begin{tablalafrox}{Diccionario del Producto 4: Software operativo. Fuente: según cita por cuenta en cada encabezado.}{tab:T14-dic-prod-4}
  {L{1.2cm} Y Y Y C{0.9cm} L{2.3cm}}
  {\cab{Código} & \cab{Paquete de trabajo} & \cab{Entregable} & \cab{Criterio de aceptación} & \cab{Resp.} & \cab{Período}}
  \multicolumn{6}{>{\columncolor{lafroxBg}\RaggedRight\arraybackslash}p{\dimexpr\textwidth-2\tabcolsep\relax}}{\textbf{Cuenta 4.1 — Bodega: recepción, inventario y preparación.} Esta cuenta reúne los módulos de recepción, inventario y preparación de bodega de la Etapa 1. {\color{lafroxGris}Fuente: (Bases Administrativas, arts. 17.2–17.6 y Formulario T-14, p. 63); (Bases Técnicas Transversales, secciones 17.2–17.6, pp. 32–34).}} \\
  4.1.1 & Recepción de mercadería con captura de lote y vencimiento (M1 Recepción) & M1 en QA. & Los requerimientos del T-12 asignados al módulo pasan sus pruebas; además, ningún producto con trazabilidad obligatoria se recibe sin lote. & DES & E1 · Meses 5 a 7 (ola A); H4. \\
  4.1.2 & Inventario por ubicación y lote, con FEFO y conteo cíclico (M2 Inventario) & M2 en QA. & Los requerimientos del T-12 asignados al módulo pasan sus pruebas; además, cada diferencia de conteo queda con su causa. & DES & E1 · Meses 5 y 6 (ola A); H4. \\
  4.1.3 & Preparación de pedidos en el turno de noche y en la cámara de congelado (M5 Preparación) & M5 en QA. & Los requerimientos del T-12 asignados al módulo pasan sus pruebas; además, se opera con guantes en el terminal de congelado, y cada faltante queda con su causa. & DES & E1 · Meses 5 a 7 (ola B); construye después de los contratos de 4.1.1/\allowbreak{}4.1.2 e integra después de su entrega (D-16). \\
  \midrule
  \multicolumn{6}{>{\columncolor{lafroxBg}\RaggedRight\arraybackslash}p{\dimexpr\textwidth-2\tabcolsep\relax}}{\textbf{Cuenta 4.2 — Calidad y cadena de frío.} Esta cuenta reúne los módulos de calidad y cadena de frío de la Etapa 1. {\color{lafroxGris}Fuente: (Bases Administrativas, arts. 17.2–17.6 y Formulario T-14, p. 63); (Bases Técnicas Transversales, secciones 17.2–17.6, pp. 32–34).}} \\
  4.2.1 & Cadena de frío, retención por SSCC y retiro sanitario (M9 Calidad y trazabilidad) & M9 y doble evaluación del TC58e en QA. & Reglas por familia, acción por defecto, exposición de andén, disposición y protocolo de ruptura; V6–V15 y V21 con retención por SSCC; retiro en menos de 2 h. & DES & E1 · Meses 5 a 7 (ola B); H4, antes del piloto. \\
  4.2.2 & Ingesta de temperatura (M9) y posición de flota (M12) & Greengrass de camión, ingesta móvil y satelital M9; posición M12 en QA. & Temperatura a cargo de M9 con IoT Core y CloudConnect/SQS; móvil y satélite siempre activos, sin pérdidas ni duplicados; M12 conserva solo posición. & DES & E1 · Meses 5 a 7 (ola B); H4, antes del piloto. \\
  \midrule
  \multicolumn{6}{>{\columncolor{lafroxBg}\RaggedRight\arraybackslash}p{\dimexpr\textwidth-2\tabcolsep\relax}}{\textbf{Cuenta 4.3 — Venta y reparto.} Esta cuenta reúne los módulos de venta y reparto de la Etapa 1. {\color{lafroxGris}Fuente: (Bases Administrativas, arts. 17.2–17.6 y Formulario T-14, p. 63); (Bases Técnicas Transversales, secciones 17.2–17.6, pp. 32–34).}} \\
  4.3.1 & Aplicación móvil de preventa sin conexión, con stock y crédito (M3 Preventa) & M3 en QA. & Los requerimientos del T-12 asignados al módulo pasan sus pruebas; además, ningún pedido se pierde ni se duplica al recuperar la señal. & DES & E1 · Meses 5 a 7 (ola A); H4. \\
  4.3.2 & Planificación de rutas corregible por el planificador (M4 Rutas) & M4 en QA. & Los requerimientos del T-12 asignados al módulo pasan sus pruebas; además, la ruta se genera en menos de 20 minutos y don Hugo la valida. & DES & E1 · Meses 7 y 8 (ola C, al terminar 4.1.1); construye después de los contratos de 4.1.3 y 4.3.1 e integra después de su entrega (D-18); H4. \\
  4.3.3 & Aplicación móvil de reparto con prueba de entrega digital (M6 Reparto) & M6 en QA. & Los requerimientos del T-12 asignados al módulo pasan sus pruebas; además, funciona sin señal, y la prueba de entrega llega al ERP el mismo día. & DES & E1 · Meses 7 y 8 (ola C, al terminar 4.3.1); construye después de los contratos de 4.1.3 y 4.3.1 e integra después de su entrega (D-18); H4. \\
  4.3.4 & Devoluciones en terreno y control de envases (M8 Devoluciones y envases) & M8 en QA. & Los requerimientos del T-12 asignados al módulo pasan sus pruebas; además, cada devolución queda con su causa. & DES & E1 · Meses 6 a 8 (ola C, al terminar 4.1.2); H4. \\
  4.3.5 & Cobro en ruta, rendición y conciliación (M7 Cobranza y rendición) & M7 en QA. & Los requerimientos del T-12 asignados al módulo pasan sus pruebas; además, cada cobro queda unido a su entrega. & DES & E1 · Meses 7 y 8 (ola C); construye después de los contratos de 4.3.3 e integra después de su entrega (D-19); H4. \\
  \midrule
  \multicolumn{6}{>{\columncolor{lafroxBg}\RaggedRight\arraybackslash}p{\dimexpr\textwidth-2\tabcolsep\relax}}{\textbf{Cuenta 4.4 — Canal moderno y portales.} Esta cuenta reúne el software de canal moderno y los portales que la Etapa 2 pone en producción en el mes 21. Se construye entre los meses 13 y 18, durante la marcha blanca de la Etapa 1, y entra al H9 (mes 17) con QA superado. {\color{lafroxGris}Fuente: (Bases Administrativas, arts. 17.2–17.6 y Formulario T-14, p. 63); (Bases Técnicas Transversales, secciones 17.2–17.6, pp. 32–34).}} \\
  4.4.1 & Pedido electrónico y aviso de despacho para las cadenas, y portal de proveedores (M11 Canal moderno) & M11 en QA. & Los requerimientos del T-12 asignados pasan sus pruebas, incluido el portal de proveedores que consulta en modo de solo lectura el estado de sus órdenes de compra (RF-12.22). & DES & E2 · Meses 14 y 15, tras la entrega del diseño del H8; integración 6.1.2 del 2 al 18 de mayo de 2028 (mes 16); entrega del expediente el 18 de mayo y acta H9 el 1 de junio de 2028 (mes 17). \\
  4.4.2 & Portal web de clientes: cuenta, saldo, entregas y documentos & Portal en QA. & Los requerimientos del T-12 asignados pasan sus pruebas; además, cada cliente ve solo sus datos, y el portal cumple las normas de accesibilidad. & DES & E2 · Meses 14 y 15, tras la entrega del diseño del H8; integración 6.1.2 del 2 al 18 de mayo de 2028 (mes 16); entrega del expediente el 18 de mayo y acta H9 el 1 de junio de 2028 (mes 17). \\
  4.4.3 & Portal web de transportistas: rutas asignadas del día siguiente & Portal en QA. & Los requerimientos del T-12 asignados pasan sus pruebas; además, cada empresa ve solo sus rutas. & DES & E2 · Meses 14 y 15, tras la entrega del diseño del H8; integración 6.1.2 del 2 al 18 de mayo de 2028 (mes 16); entrega del expediente el 18 de mayo y acta H9 el 1 de junio de 2028 (mes 17). \\
  \midrule
  \multicolumn{6}{>{\columncolor{lafroxBg}\RaggedRight\arraybackslash}p{\dimexpr\textwidth-2\tabcolsep\relax}}{\textbf{Cuenta 4.5 — Analítica (M10).} Esta cuenta reúne las capacidades de analítica M10 de las dos etapas: indicadores operacionales y OTIF diario para la Etapa 1, y costo de servir por cliente y por entrega para la Etapa 2. {\color{lafroxGris}Fuente: (Bases Administrativas, arts. 17.2–17.6 y Formulario T-14, p. 63); (Bases Técnicas Transversales, secciones 17.2–17.6, pp. 32–34).}} \\
  4.5.1 & Indicadores operacionales y OTIF diario (M10 Analítica) & Indicadores publicados. & Los requerimientos del T-12 asignados al módulo pasan sus pruebas; además, operaciones reproduce el cálculo sobre una muestra. & DAT & E1 · Meses 7 y 8 (ola C, al terminar 4.2.1); H4. \\
  4.5.2 & Costo de servir por cliente y por entrega (M10 Analítica) & Costo de servir publicado. & Los requerimientos del T-12 asignados pasan sus pruebas; además, finanzas reproduce el costo de una muestra de entregas. & DAT & E2 · Meses 14 y 15, tras la entrega del diseño del H8; integración 6.1.2 del 2 al 18 de mayo de 2028 (mes 16); entrega del expediente el 18 de mayo y acta H9 el 1 de junio de 2028 (mes 17). \\
\end{tablalafrox}
